%!TEX root = ./main.tex

\section[DFRC]{From Two Boxes to One}

% SECTION TIMING: 12:00--22:00 (4 frames).
% This section earns the rest of the deck. If the class cannot say what delay and Doppler
% measure, the waveform section will not land. Do not cut the "what an echo tells you" frame.

\begin{frame}{Two boxes became one box}
  \vspace{0.3em}
  \centering
  \begin{tikzpicture}[scale=0.88,transform shape]
    % separate
    \node[font=\small\bfseries,anchor=south] at (-2.3,1.55) {Separate systems};
    \node[netnode,minimum width=1.9cm] (c1) at (-3.4,0.55) {Comms};
    \node[corenode,minimum width=1.9cm] (r1) at (-3.4,-0.75) {Radar};
    \node[font=\scriptsize,align=left,anchor=west] at (-2.35,0.55) {own band\\own antenna};
    \node[font=\scriptsize,align=left,anchor=west] at (-2.35,-0.75) {own band\\own antenna};
    % divider
    \draw[softgray,dashed] (0.45,1.6) -- (0.45,-1.6);
    % shared
    \node[font=\small\bfseries,anchor=south] at (3.6,1.55) {DFRC};
    \node[netnode,minimum width=2.6cm,minimum height=1.6cm,draw=coreorange,
      fill=orange!5] (d) at (2.9,-0.1) {One transmit\\chain};
    \node[font=\scriptsize,anchor=west] at (4.35,0.35) {\textcolor{ranblue}{bits out}};
    \node[font=\scriptsize,anchor=west] at (4.35,-0.55) {\textcolor{coreorange}{echo in}};
  \end{tikzpicture}

  \vspace{0.8em}
  \begin{columns}[T,onlytextwidth]
    \column{0.47\textwidth}
      {\small\textbf{\textcolor{softgreen}{What it buys}}\par}
      \vspace{0.3em}
      {\small One spectrum licence, one amplifier, one aperture, one power budget.\par}
    \column{0.47\textwidth}
      {\small\textbf{\textcolor{coreorange}{What it costs}}\par}
      \vspace{0.3em}
      {\small A single waveform now has to satisfy two specifications at once.\par}
  \end{columns}

  \glossline{DFRC = Dual-Function Radar Communication}
\end{frame}

\begin{frame}{DFRC is one box. ISAC is the question.}
  \vspace{0.4em}
  \centering
  \begin{tikzpicture}[scale=0.95,transform shape]
    \node[draw=coreorange,fill=orange!4,rounded corners=6pt,
      minimum width=8.2cm,minimum height=2.9cm] (isac) at (0,0) {};
    \node[anchor=north west,font=\bfseries\small,text=coreorange] at (-4.0,1.32) {ISAC};
    \node[draw=ranblue,fill=blue!6,rounded corners=4pt,
      minimum width=3.4cm,minimum height=1.5cm,align=center,font=\bfseries\small]
      (dfrc) at (-1.9,-0.25) {DFRC\\[-2pt]{\scriptsize one shared box}};
    \node[align=center,font=\scriptsize,text width=3.2cm] at (1.9,-0.25)
      {separate devices,\\separate sites,\\one sensing picture};
  \end{tikzpicture}

  \vspace{0.7em}
  \begin{columns}[T,onlytextwidth]
    \column{0.47\textwidth}
      {\small\textbf{\textcolor{ranblue}{DFRC}} shares waveforms, antennas, and
       processing chains inside one radio.\par}
    \column{0.47\textwidth}
      {\small\textbf{\textcolor{coreorange}{ISAC}} also lets one node transmit
       and a different node listen.\par}
  \end{columns}

  \vspace{0.5em}
  \takeaway{Every DFRC system is ISAC.\\Not every ISAC system fits in one box.}
\end{frame}

% Teaching note (2:00): this frame is not in the source talk. It is here because the
% waveform section is unreadable without it. Draw the three arrows on the board as you go.
\begin{frame}{What an echo actually tells you}
  \vspace{0.4em}
  \begin{columns}[T,onlytextwidth]
    \column{0.44\textwidth}
      \centering
      \begin{tikzpicture}[scale=0.85,transform shape]
        \node[netnode,minimum width=1.5cm] (bs) at (0,0) {BS};
        \node[corenode,minimum width=1.5cm] (tg) at (3.6,0) {Target};
        \draw[flow,draw=coreorange,very thick] (0.78,0.28) -- (2.82,0.28);
        \draw[draw=coreorange,very thick,dashed,-{Latex[length=2.2mm]}]
          (2.82,-0.28) -- (0.78,-0.28);
        \node[font=\scriptsize,anchor=south] at (1.8,0.35) {transmit};
        \node[font=\scriptsize,anchor=north] at (1.8,-0.35) {echo};
      \end{tikzpicture}
    \column{0.52\textwidth}
      \small
      \begin{itemize}\itemsep6pt
        \item \textbf{Delay} $\rightarrow$ range, $R = c\tau/2$
        \item \textbf{Doppler} $\rightarrow$ closing speed
        \item \textbf{Angle} $\rightarrow$ direction
      \end{itemize}
  \end{columns}

  \vspace{0.6em}
  \qa{Which knob buys finer range, and which buys finer speed?}
     {\textbf{Bandwidth} buys range resolution --- a wider signal has a sharper
      peak in time. \textbf{Observation time} buys velocity resolution --- you
      must watch long enough for the phase to drift.}
\end{frame}

\begin{frame}{The four hard parts, named up front}
  \vspace{0.5em}
  \begin{columns}[T,onlytextwidth]
    \column{0.47\textwidth}
      \large\textbf{\textcolor{ranblue}{Techniques}}

      \vspace{0.5em}
      \small
      \begin{itemize}\itemsep4pt
        \item Waveform design
        \item MIMO, beamforming
        \item Shared RF hardware
        \item Coding and allocation
      \end{itemize}
    \column{0.47\textwidth}
      \large\textbf{\textcolor{coreorange}{Challenges}}

      \vspace{0.5em}
      \small
      \begin{itemize}\itemsep4pt
        \item Tradeoff: unavoidable
        \item Interference: mutual
        \item Complexity: non-convex
        \item Regulation: who owns it
      \end{itemize}
  \end{columns}

  \glossline{MIMO = Multiple Input Multiple Output \quad RF = Radio Frequency}

  \vspace{0.4em}
  \takeaway{Three of the four are engineering.\\The fourth is a committee.}
\end{frame}
